\documentclass[10pt,a4paper]{article}
\usepackage[T1]{fontenc}
\usepackage[width=20cm, height=30cm, left=20mm, right=20mm, top=10mm, bottom=10mm]{geometry}
\usepackage{graphicx}
\usepackage{mathtools}
\usepackage{amssymb}
\usepackage{ctex}
\usepackage{nameref}

\newcommand{\R}{\mathbb{R}}
\newcommand{\diff}{\mathop{}\!\mathrm{d}}
\newcommand{\geh}[1]{\sqrt{#1}}
\newcommand{\hs}[2]{#1(#2)}
\newcommand{\sigm}[2]{\sum^{#2}_{i=#1}}
\newcommand{\xiab}[1]{_{#1}}
\newcommand{\p}[2]{#1^{'}(#2)}
\newcommand{\pp}[2]{#1^{"}(#2)}
\newcommand{\so}{\therefore}
\newcommand{\be}{\because}

\begin{document}
$	\hs{h}{x}=\geh{x+45}-\geh{x+5}$
求 $[\sigm{1}{n}x\xiab{i}]$

\subsection*{g(x)性质}

$\p{g}{x}=\frac{1}{2\geh{x+45}}-\frac{1}{2\geh{x+5}}$
$\pp{g}{x}=- \frac{1}{4(x+45)^{3/2}} + \frac{1}{4(x+5)^{3/2}}>0$

$x>0$时，g(x)\textuparrow g(x)下凸

\subsection*{背景}
%插入图片1
由第一问得 $3.9<x\xiab{2n}<4.0
 ,  4.0<x\xiab{2n=1}<4.1$
发现$|x\xiab{n}-4|\textdownarrow$ 且 $\textrightarrow $0
\subsection*{构造直线}
$\p{g}{4}=-\frac{2}{21}$

构造3条直线：
$$\hs{a}{x}=-\frac{2}{21}(x-4)+4$$
$$\hs{a}{x}=-\frac{1}{4}(x-4)+4$$
$$\hs{a}{x}=-k(x-4)+4$$
%插入图片2
易证：$$0<x<4,a(x)<g(x)<b(x)   (01)$$
$$4<x<5,a(x)<g(x)<c(x)           (02)$$
\subsection*{放缩}
%插入图片3

$\because 0<x\xiab{2k}<4$ 且$ (01)$

$\therefore a(x_{2k})<g(x_{2k})<b(x_{2k})$

$\because a(x)\textdownarrow$

$\so a(x_{2k})>a(4)=4$

$\be g(x_{2k})=x_{2k+1}$

$\so 4<a(x\xiab{2k})<x_{2k+1}<b(x\xiab{2k})$

即$$-\frac{2}{21}(x_{2k}-4)+4<x_{2k+1}<-\frac{1}{4}(x_{2k}-4)+4     (1)$$

令$k=1$ $$-\frac{2}{21}(x_{2}-4)<x_{3}-4<-\frac{1}{4}(x_{2}-4)    (7)$$


%插入图片4
$
\be 0<X<4,c(x)>g(x)>a(x)$

$\so c(a(x_{2k}))>g(a(x_{2k})) ,  g(x_{2k+1})>a(x\xiab{2k=1})$

$\be c,g,a\textdownarrow$

$\so g(a(x_{2k}))>
g(x_{2k+1}),a(x\xiab{2k+1})>a(b(x\xiab{2k}))
$

$\so c(a(x_{2k}))>g(x_{2k+1})>a(b(x\xiab{2k}))$

即$$ -2/21[-1/4(x\xiab{2k}-4)]+4
<x\xiab{2k+2}<
-k[-2/21(x\xiab{2k-4})]+4   
$$
$$\frac{1}{42}(x_{2k}-4)<x\xiab{2k+2}-4<\frac{2}{21}k(x\xiab{2k-4})(2)$$

$$\frac{1}{1-1/42}(x\xiab{2}-4)
<\frac{1-(\frac{1}{42})^{n}}{1-\frac{1}{42}}(x\xiab{2}-4)
<\sigm{1}{n} x\xiab{2k}-4
<\frac{1-(\frac{2k}{21})^{n}}{1-\frac{2k}{21}}(x\xiab{2}-4)
<(x\xiab{2}-4)
$$
即$$\frac{1}{1-1/42}(x\xiab{2}-4)<
\sigm{1}{n} x\xiab{2k}-4
<(x\xiab{2}-4)                       (3)
$$

类似的，可得
$$-\frac{2}{21}(x_{2k-1}-4)+4<x_{2k}<-k(x_{2k-1}-4)+4     ()$$

令$k=1$ $$-\frac{2}{21}(x_{1}-4)<x_{2}-4<-k(x_{1}-4)    (9)$$






$$\frac{1}{1-1/42}(x\xiab{3}-4)
>\frac{1-(\frac{1}{42})^{n}}{1-\frac{1}{42}}(x\xiab{3}-4)
>\sigm{1}{n} x\xiab{2k+1}-4
>\frac{1-(\frac{2k}{21})^{n}}{1-\frac{2k}{21}}(x\xiab{3}-4)
>(x\xiab{3}-4)
$$
即$$\frac{1}{1-1/42}(x\xiab{3}-4)>
\sigm{1}{n} x\xiab{2k+1}-4
>(x\xiab{3}-4)                       (6)
$$


由$(3)(6)$
$$(x\xiab{3}-4)+\frac{1}{1-1/42}(x\xiab{2}-4)<
\sigm{2}{n} x\xiab{i}-4<
\frac{1}{1-1/42}(x\xiab{3}-4)+(x\xiab{2}-4)  
$$
由$(7)$
$$(-\frac{2}{21}+\frac{1}{1-\frac{1}{42}})(x\xiab{2}-4)
<\sigm{2}{n} x\xiab{i}-4<
(\frac{1}{1-\frac{1}{42}}(-\frac{1}{4})+1)(x\xiab{2}-4)$$

$\be  0<-\frac{2}{21}+\frac{1}{1-\frac{1}{42}}<1$且$(x\xiab{2}-4)<0$


$\so (-\frac{2}{21}+\frac{1}{1-\frac{1}{42}})(x\xiab{2}-4)>
(x\xiab{2}-4)$


$\be \frac{1}{1-\frac{1}{42}}(-\frac{1}{4})+1>0$且$(x\xiab{2}-4)<0$


$\so (\frac{1}{1-\frac{1}{42}}(-\frac{1}{4})+1)(x\xiab{2}-4)<0$


得$$
(x\xiab{2}-4)<
\sigm{2}{n} x\xiab{i}-4
<0$$

由$(9)$
$$
-\frac{2}{21}(x\xiab{1}-4)<
\sigm{2}{n} x\xiab{i}-4
<0$$
即
$$
-\frac{2}{21}<
\sigm{2}{n} x\xiab{i}-4
<0$$

$\so $ $$
-\frac{2}{21}<
(\sigm{2}{n} x\xiab{i})-4(n-1)
<0$$

$$
-\frac{2}{21}+4(n-1)<
(\sigm{2}{n} x\xiab{i})
<4(n-1)$$

$$
1-\frac{2}{21}+4n<
(\sigm{2}{n} x\xiab{i})+5
<4n+1$$

$$
1-\frac{2}{21}+4n<
\sigm{1}{n} x\xiab{i}
<4n+1$$

$\so $
$$[\sigm{1}{n}x\xiab{i}]= 4n$$
\end{document}